\setcounter{section}{2}

\Section{Предел числовой последовательности. Единственность предела.
Бесконечно малые последовательности и их свойства. Свойства пределов,
связанные с неравенствами. Арифметические операции со сходящимися
последовательностями. Теорема Вейерштрасса о пределе монотонной
ограниченной последовательности. Число \(e\). Бесконечно малые и бесконечно
большие последовательности и их свойства}

\Subsection{Бесконечно малые и бесконечно большие последовательности,
их свойства}

\begin{definition}
    Последовательность \(\{x_{n}\}\) называют \textit{бесконечно малой},
    если для любого \(\varepsilon>0\) существует номер \(N\) такой, что
    при \(n\geq N\) выполняется неравенство \(|x_{n}|<\varepsilon\).
    \[[\{x_{n}\}\text{ - бесконечно малая}]\overset{def}{=}
    [\forall\,\epsilon>0~~\exists\,N\in\N:\forall n\in\N,~n\geq N\longmapsto|x_{n}|<\epsilon]\]
    Далее бесконечно малую последовательность будем обозначать БМП.
\end{definition}

\begin{definition}
    Последовательность \(\{x_{n}\}\) называют \textit{бесконечно большой},
    если для любого \(\gamma>0\) существует номер \(N\) такой, что при \(n\geq N\)
    выполняется неравенство \(|x_{n}|\geq\gamma\).
    \[[\{x_{n}\}\text{ - бесконечно большая}]\overset{def}{=}
    [\forall\,\gamma>0~~\exists N\in\N:\forall n\in\N,~n\geq N\longmapsto|x_{n}|\geq\gamma]\]
    Далее бесконечно большую
    последовательность будем обозначать ББП.
\end{definition}

\vspace{3mm}
Добавим сюда же определение, которое не задействовано конкретно в этой теме,
но используется в будущем.

\begin{definition}
    Последовательность \(\{x_{n}\}\) называется \textit{неограниченной сверху},
    если для любого действительного числа \(\alpha\) найдется такой номер \(N\),
    что выполняется неравенство \(x_{N}>\alpha\).
    \[\forall\,\alpha\in\R~~\exists\,N\in\N:~~x_{N}>\alpha.\]  
\end{definition}
Неограниченность снизу определяется аналогично. 
\begin{anote}
    Почувствуйте разницу между неограниченными и бесконечно большими последовательностями:
    у первых \textit{хотя бы} один член больше нужного числа, а у ББП \textit{все}
    члены после некоторого номера по модулю больше нужного числа. \\
    \indent Например, последовательность, которая на четных номерах равна \(0\),
    а на нечетных равна \(n\), является неограниченной сверху,
    но не является бесконечно большой.
\end{anote}

\vspace{3mm}

\begin{theorem}
    Сумма и разность двух БМП являются БМП.

    \vspace{2mm}

    \proofname~Пусть \(\{x_{n}\}\) и \(\{y_{n}\}\) --- БМП. Для заданного
    \(\varepsilon>0\) выберем номера \(N_{1}\) и \(N_{2}\), начиная с
    которых \(|x_{n}|<\frac{\varepsilon}{2}\) и
    \(|y_{n}|<\frac{\varepsilon}{2}\) соответственно. При
    \(n\geq \max(N_{1},N_{2})\) имеем
    \[|x_{n}\pm y_{n}|\leq |x_{n}|+|y_{n}|<\varepsilon.\]
    Следовательно, сумма и разность являются БМП.~\qedsymbol
\end{theorem}

\begin{theorem}
    Произведение ограниченной последовательности и БМП является БМП.

    \vspace{2mm}

    \proofname~Пусть \(|x_{n}|\leq C\) для любого \(n\), а
    \(\{y_{n}\}\) --- БМП. Если \(C=0\), то произведение равно нулю.
    Если \(C>0\), то для заданного \(\varepsilon>0\) начиная с некоторого
    номера \(N\) выполняется
    \[|y_{n}| < \frac{\epsilon}{C}\]
    \[|x_{n}y_{n}|\leq C|y_{n}|<\varepsilon.\]
    Значит, произведение является БМП.~\qedsymbol
\end{theorem}

\begin{theorem}
    Любая БМП ограничена.

    \vspace{2mm}

    \proofname~Для \(\varepsilon=1\) найдётся номер \(N\), после которого
    \(|x_{n}|<1\). Первые \(N\) членов образуют конечное множество и
    потому ограничены. Значит, ограничена вся последовательность.~\qedsymbol
\end{theorem}

\begin{corollary}
    Произведение конечного количества БМП является БМП.
\end{corollary}

\begin{theorem}
    Если все члены БМП, начиная с некоторого номера, равны одному и тому же
    числу \(\gamma\), то \(\gamma=0\).

    \vspace{2mm}

    \proofname~Предположим, что \(\gamma\neq 0\), и возьмём
    \(\varepsilon=\frac{|\gamma|}{2}\). Тогда существует такой номер \(N\),
    что для любых \(n\geq N\) должно выполняться \(|x_{n}|<\varepsilon\), но
    \(|x_{n}|=|\gamma|>\varepsilon\). Получено противоречие.~\qedsymbol
\end{theorem}

\begin{theorem}
    Если \(\{x_{n}\}\) --- ББП, то, начиная с некоторого номера, определена
    последовательность \(\left\{\frac{1}{x_{n}}\right\}\), являющаяся БМП.

    \vspace{2mm}

    \proofname~Так как \(\{x_{n}\}\) --- ББП, то начиная с некоторого
    номера \(|x_{n}|\geq1\), поэтому \(x_{n}\neq 0\) и обратная
    последовательность определена. Для любого \(\varepsilon>0\) начиная
    с некоторого номера \(|x_{n}|\geq\frac{2}{\varepsilon}\), откуда
    \[\left|\frac{1}{x_{n}}\right|\leq\frac{\varepsilon}{2}<\varepsilon.\]
    Следовательно, обратная последовательность является БМП.~\qedsymbol
\end{theorem}

\begin{theorem}
    Если \(\{x_{n}\}\) --- БМП и каждый её член не равен нулю, то
    последовательность \(\left\{\frac{1}{x_{n}}\right\}\) --- ББП.

    \vspace{2mm}

    \proofname~Для любого \(\gamma>0\) из определения БМП следует, что начиная
    с некоторого номера \(|x_{n}|<\frac{1}{\gamma}\). Поскольку \(x_{n}\neq 0\),
    получаем
    \[\left|\frac{1}{x_{n}}\right|\geq\gamma.\]
    Значит, обратная последовательность является ББП.~\qedsymbol
\end{theorem}

\begin{theorem}
    Любая подпоследовательность ББП является ББП.

    \vspace{2mm}

    \textit{Примечание: смотри доказательство в Теореме 4.3}
\end{theorem}

\Subsection{Предел числовой последовательности. Единственность предела}

\begin{definition}
    Число \(a\in\R\) называют \textit{пределом числовой последовательности}
    \(\{x_{n}\}\), а последовательность --- \textit{сходящейся к числу}
    \(a\), если для любого \(\varepsilon>0\) существует номер \(N\) такой,
    что при \(n>N\) выполняется неравенство \(|x_{n}-a|<\varepsilon\).
    Обозначают \({\lim\limits_{n\to\infty}x_{n}=a}\).
\end{definition}

\begin{note}
    Любую сходящуюся последовательность можно представить в виде суммы
    её предела и бесконечно малой последовательности:
    \[x_{n}=a+a_{n},\qquad \lim\limits_{n\to\infty}x_{n}=a.\]
\end{note}

\begin{theorem}
    Любая сходящаяся числовая последовательность имеет единственный предел.

    \vspace{2mm}

    \proofname~Предположим, что последовательность \(\{x_{n}\}\) имеет
    два предела \(a\) и \(b\). По предыдущему замечанию существуют БМП
    \(\{a_{n}\}\) и \(\{b_{n}\}\) такие, что
    \[a+a_{n}=x_{n}=b+b_{n}.\]
    Тогда \(a-b=b_{n}-a_{n}\) для любого \(n\). По теореме о сумме и
    разности БМП последовательность \(\{b_{n}-a_{n}\}\) является БМП.
    Следовательно, постоянная последовательность со значением \(a-b\)
    является БМП. По теореме о постоянных членах БМП получаем \(a-b=0\),
    то есть \(a=b\). Поэтому предел единственный.~\qedsymbol
\end{theorem}

\Subsection{Арифметические операции со сходящимися последовательностями.
Свойства пределов, связанные с неравенствами}

Сначала рассмотрим арифметические операции.

\begin{theorem}
    Сумма (разность) сходящихся числовых последовательностей является
    сходящейся числовой последовательностью. Предел суммы (разности)
    равен сумме (разности) пределов этих последовательностей.

    \vspace{2mm}

    \proofname~Пусть \(x_{n}=a+a_{n}\) и \(y_{n}=b+b_{n}\), где \(a_{n}\) и \(b_{n}\) - БМП.
    По теореме о сумме и разности БМП последовательность
    \[\gamma_{n}=(x_{n}\pm y_{n})-(a\pm b)=(x_{n}- a)\pm(y_{n}- b)=a_{n}-b_{n}\]
    является БМП. Но тогда отсюда следует:
    \[\lim_{ n \to \infty }(x_{n}-y_{n}) = 
    \lim_{ n \to \infty }(\gamma_{n}+a\pm b) =
    \lim_{ n \to \infty } \gamma_{n}+(a\pm b) = a\pm b.\]
    Значит,
    \(\lim\limits_{n\to\infty}(x_{n}\pm y_{n})=a\pm b\).
    ~\qedsymbol
\end{theorem}

\begin{theorem}
    Произведение конечного числа сходящихся числовых последовательностей
    является сходящейся числовой последовательностью, а его предел
    равен произведению пределов этих последовательностей.

    \vspace{2mm}

    \proofname~Докажем утверждение для двух последовательностей.
    Пусть \(x_{n}=a+a_{n}\) и \(y_{n}=b+b_{n}\), где \(a_{n}\) и \(b_{n}\) - БМП. 
    Последовательность \(\{x_{n}\}\) ограничена, так как ограничены \(a\) и \(\{a_{n}\}\).
    Зададим следующую последовательность:
    \[\gamma_{n}=x_{n}y_{n}-ab=x_{n}b_{n}+b(x_{n}-a)=x_{n}b_{n}+a_{n}b\]
    По теоремам 3.1 и 3.2 получаем, что \(\gamma_{n}\) - БМП. Следовательно,
    \[\lim_{ n \to \infty }x_{n}y_{n} =
    \lim_{ n \to \infty }(\gamma_{n}+ab) =
    \lim_{ n \to \infty } \gamma_{n} + ab = ab\]
    Для общего случая требуется лишь применить доказанное нужное количество раз.~\qedsymbol
\end{theorem}

\begin{theorem}
    Если \(\lim\limits_{n\to\infty}x_{n}=\alpha\) и
    \(\lim\limits_{n\to\infty}y_{n}=\beta\neq 0\), то начиная
    с некоторого номера определена последовательность
    \(\left\{\frac{x_{n}}{y_{n}}\right\}\), причём
    \[\lim\limits_{n\to\infty}\frac{x_{n}}{y_{n}}=\frac{\alpha}{\beta}.\]

    \vspace{2mm}

    \proofname~Так как
    \(\lim\limits_{n\to\infty}y_{n}=\beta\neq 0\), то имеем:
    \[\epsilon=\frac{|\beta|}{2}~~\exists\,N_{1}\in\N:
    \forall n\in\N,~n\geq N_{1}\longmapsto|y_{n}-\beta|<\epsilon=\frac{|\beta|}{2},\]
    \[\forall\,\epsilon>0~~\exists\,N_{2}\in\N:
    \forall n\in N,~n\geq N_{2}\longmapsto|y_{n}-\beta|<\frac{\epsilon|\beta|^2}{2}.\]
    Выберем \(N=\max(N_{1},N_{2})\), чтобы выполнялись оба утверждения.
    Первое нам нужно для того, чтобы мы точно не разделили на ноль.
    Оба утверждения нужны для следующих преобразований 
    (во второй строке используем сначала первое, потом второе):
    \[|y_{n}-\beta|<\frac{|\beta|}{2}~\Rightarrow~|y_{n}|>\frac{|\beta|}{2}\]
    \[\left|\frac{1}{y_{n}}-\frac{1}{\beta}\right|
    =\frac{|y_{n}-\beta|}{|y_{n}||\beta|}
    <\frac{2|y_{n}-\beta|}{|\beta|^{2}} < \epsilon.\]
    Тогда по определению \(\lim\limits_{n\to\infty}\frac{1}{y_{n}}=\frac{1}{\beta}\).
    Далее по теореме о произведении сходящихся последовательностей получаем
    \[\lim\limits_{n\to\infty}\frac{x_{n}}{y_{n}}=
    \lim_{ n \to \infty } \frac{1}{y_{n}}x_{n}=\frac{1}{\beta}\alpha=\frac{\alpha}{\beta}\]
    Что и требовалось доказать.~\qedsymbol
\end{theorem}

Теперь переходим к свойствам сходящихся последовательностей, связанным с неравенствами.

\begin{theorem}
    Если \(\lim\limits_{n\to\infty}x_{n}=\alpha\) и начиная с некоторого
    номера \(x_{n}\geq\beta\),
    то \(\alpha\geq\beta\).

    \vspace{2mm}

    \proofname~Предположим, что \(\alpha<\beta\), и положим
    \(\varepsilon=\frac{\beta-\alpha}{2}\). Начиная с некоторого
    номера \(N\) одновременно выполняются \(|x_{n}-\alpha|<\varepsilon\)
    и \(x_{n}\geq\beta\). Первое неравенство даёт
    \(x_{n}<\alpha+\varepsilon<\beta\), что противоречит второму.
    Следовательно, \(\alpha\geq\beta\).~\qedsymbol
\end{theorem}

\begin{theorem}
    Если начиная с некоторого номера выполняется \(x_{n}\leq y_{n}\),
    где \({\lim\limits_{n\to\infty}x_{n}=\alpha}\) и
    \(\lim\limits_{n\to\infty}y_{n}=\beta\), то \(\alpha\leq\beta\).

    \vspace{2mm}

    \proofname~Предположим, что \(\alpha>\beta\), и положим
    \(\varepsilon=\frac{\alpha-\beta}{3}\). Начиная с некоторого номера \(N\) 
    для всех номеров \(n\geq N\) имеем
    \[x_{n}>\alpha-\varepsilon>\beta+\varepsilon>y_{n},\]
    что противоречит условию \(x_{n}\leq y_{n}\). Значит,
    \(\alpha\leq\beta\).~\qedsymbol
\end{theorem}

\begin{theorem}
    Если все члены сходящейся числовой последовательности \(\{x_{n}\}\)
    принадлежат отрезку \([a,b]\), то и её предел принадлежит этому
    отрезку.

    \vspace{2mm}

    \proofname~Пусть \(\lim\limits_{n\to\infty}x_{n}=c\). Из \(x_{n}\geq a\) и
    Tеоремы 3.12 следует \(c\geq a\). Кроме того, \(-x_{n}\geq-b\) и
    \(\lim\limits_{n\to\infty}(-x_{n})=-c\), поэтому по той же теореме \(-c\geq-b\),
    то есть \(c\leq b\). Следовательно, \(c\in[a,b]\).~\qedsymbol
\end{theorem}

\Subsection{Теорема Вейерштрасса о пределе монотонной ограниченной
последовательности}

\begin{theorem}[Теорема Вейерштрасса]
    Если последовательность не убывает и ограничена сверху, то она
    сходится, причём \(\lim\limits_{ n \to \infty }{x_{n}} = \underset{n}{\sup}~x_{n}\). 
    Если последовательность не возрастает и ограничена снизу, то
    она сходится, причём \(\lim\limits_{ n \to \infty }{x_{n}} = \underset{n}{\inf}~x_{n}.\)

    \vspace{2mm}

    \proofname~Докажем утверждение для неубывающей последовательности
    \(\{x_{n}\}\), ограниченной сверху. Обозначим через \(a\) точную
    верхнюю грань множества её членов:
    \[a=\underset{n}{\sup}~x_{n}\]
    Для любого \(\varepsilon>0\) число \(a-\varepsilon\) не является
    верхней гранью этого множества. Поэтому найдётся номер \(N\) такой,
    что \(x_{N}>a-\varepsilon\). Так как последовательность не убывает,
    то при \(n\geq N\)
    \[a-\varepsilon<x_{N}\leq x_{n}\leq a.\]
    Следовательно, \(|x_{n}-a|<\varepsilon\), то есть
    \(\lim\limits_{n\to\infty}x_{n}=a\).

    Случай не возрастающей последовательности, ограниченной снизу,
    доказывается аналогично.~\qedsymbol
\end{theorem}

\Subsection{Число \(e\)}

Рассмотрим последовательность
\(\{x_{n}=\left(1+\frac{1}{n}\right)^{n}\}\).

\begin{theorem}
    Последовательность \(\{x_{n}\}\) сходится.

    \vspace{2mm}

    \proofname~При \(n\geq3\) по формуле бинома Ньютона имеем
    \[x_{n}=1+1+\frac{1}{2!}\left(1-\frac{1}{n}\right)
    +\frac{1}{3!}\left(1-\frac{1}{n}\right)
    \left(1-\frac{2}{n}\right)+\dots+
    \frac{1}{n!}\left(1-\frac{1}{n}\right)\dots
    \left(1-\frac{n-1}{n}\right).\]
    Первые два слагаемых дают \(2\), а остальные положительны, поэтому
    \(x_{n}>2\) при \(n\geq2\).

    Все произведения скобок с \(n\) в записи не превосходят \(1\). Для \(k\geq1\)
    имеем \(k!\geq2^{k-1}\), поэтому слагаемые сверху ограничены
    соответствующими дробями со степенями двойки:
    \[x_{n}<1+1+\frac{1}{2}+\frac{1}{4}+\frac{1}{8}+\dots=3.\]
    \indent Заметим, что при увеличении \(n\) каждое новое слагаемое положительно.
    Значит, последовательность не убывает и ограничена сверху. 
    По теореме Вейерштрасса она сходится - что и требовалось доказать.
    ~\qedsymbol
\end{theorem}

\begin{definition}
    Число \(e\) определяют как предел последовательности
    \(\left\{\left(1+\frac{1}{n}\right)^{n}\right\}\):
    \[e\overset{def}{=}\lim\limits_{n\to\infty}
    \left(1+\frac{1}{n}\right)^{n}=2{,}71828182\ldots\]
\end{definition}
